\section{Introduction and Motivation}
\label{sec:intro}

Coastal and insular ecosystems within tropical marine sanctuaries, such as the Tayrona National Natural Park (TNNP) in the Colombian Caribbean, face unprecedented stress from the influx of terrigenous sediment plumes, microplastics, and particulate organic matter \cite{bayraktarov2014physical, montoya2018seasonal}. While the macroscopic hydrodynamic forcing fields---such as the seasonal Guajira upwelling system---are widely monitored via regional implementations of continuum oceanic models (e.g., ROMS, Delft3D), a critical computational bottleneck remains at the benthic boundary layer. Traditional environmental engineering frameworks standardly jump directly to macro--scale continuum approximations, treating particulate trapping, deposition, and benthic clearance via highly simplified, empirical sink terms. Consequently, these macro--models inherently obscure the fundamental physical mechanisms governing particle--fluid and particle--particle interactions at the micro--scale.

At the coral canopy and sediment--water interfaces, local fluid transport transitions strictly into low Reynolds number regimes ($Re \ll 1$), where viscous forces and hydrodynamic lubrication dominate. Recent global physical investigations have confirmed that complex benthic structures act as highly efficient, geometry--dependent traps for colloidal and particulate pollutants through explicit hydrodynamic mechanisms, including direct boundary interception and wake--vortex stagnation \cite{mendrik2024microplastic}. In systems with high particulate volume fractions, such as the clay--colloid flocculation fronts or the high--permeability carbonate sand matrices characteristic of the Tayrona bays, particles cannot be modeled as isolated, passive tracers; instead, their time--dependent trajectories are heavily coupled via multi--body hydrodynamic interactions \cite{brady1988stokesian}. 

To bridge this fundamental disconnect between discrete micro--mechanics and macro--scale environmental forecasting, this research introduces a rigorous, multi--scale computational framework. We utilize highly optimized, open--source Stokesian Dynamics (SD) formulations \cite{brady1988stokesian, durlofsky1987dynamic, sierou2001accelerated, banchio2003accelerated} to explicitly simulate the many--body hydrodynamic interactions and lubrication forces governing colloidal suspensions within realistic, low--Reynolds complex geometries. To scale these discrete microstructural insights into metrics usable by regional engineers, we are proposing to develop custom, high--order Finite--Difference (FD) numerical solvers to evaluate the continuum advection--diffusion--reaction (ADR) equations. By mathematically matching the discrete particle--tracking distributions to the continuous scalar fields, we upscale the underlying low--level physics to derive mathematically rigorous, effective hydrodynamic dispersion tensors ($D_{\text{eff}}$) and dynamic trapping rates ($R$). Ultimately, this basic research paradigm eliminates the \emph{ad hoc} parameterization  of regional environmental modeling, providing the precise physical laws required to anchor Caribbean coastal conservation to deterministic fluid mechanics.
